%\needspace{4\baselineskip}
\color{uniblue}{\subsection{Максимальная токовая защита трансформатора (МТЗ)}}
\color{black}

\begin{enumerate}[label=\arabic{section}.\arabic{subsection}.\arabic*, labelsep=4pt, leftmargin=0pt, itemindent=57pt, itemsep=0pt, parsep=5pt] % если что можно крутить itemsep
%%%%%%%%%%%%%%%%%%%%%%%%%%%%%%%%%%%%%%%%%%%%%%%%%%%%%%%%%%%  ОБЩИЕ СВЕДЕНИЯ %%%%%%%%%%%%%%%%%%%%%%%%%%%%%%%%%%%%%%%%%%%%%%%%

\item Общие сведения

\begin{enumerate}[label=\arabic{section}.\arabic{subsection}.\arabic{enumi}.\arabic*, labelsep=4pt, leftmargin=0em, itemindent=65pt, parsep=0pt]

\item
Максимальная токовая защита <<МТЗ>> применяется для защиты обмоток трансформатора от сверхтоков внешних КЗ при повреждениях на присоединениях прилегающей сети, а также для резервирования основных защит трансформатора и подключается к трансформаторам тока стороны ВН трансформатора. Функция <<МТЗ>> имеет возможность работы с контролем напряжения стороны НН трансформатора.
\item
В составе функционального блока <<МТЗ>> предусмотрены следующие функции:
\begin{itemize}
\item три ступени максимальной токовой защиты;
\item комбинированный пусковой орган напряжения;
\item блокировка при броске намагничивающего тока;
\item блокировка логической защиты шин.
\end{itemize}

\item
В целях предотвращения действия ступеней МТЗ при коротких замыканиях в сетях 6-35 кВ (в случаях, когда нейтраль трансформатора заземлена) предусматривается возможность  исключения токов нулевой последовательности посредством переключения измеряемых токов с фазных на линейные, для этого программный переключатель SGF1 <<Сборка\_ток\_цепей>> переводится в состояние <<Треугольник>>.

\item
Функциональный блок <<МТЗ>> может быть полностью выведен из работы оперативно путем активации сигнала <<ОВ МТЗ>>.

\end{enumerate}

%%%%%%%%%%%%%%%%%%%%%%%%%%%%%%%%%%%%%%%%%%%%%%%%%%%%%%%%%  СТУПЕНЬ МТЗ %%%%%%%%%%%%%%%%%%%%%%%%%%%%%%%%%%%%%%%%%%%%%%%%

%\color{uniblue}{\subsubsection{Ступени максимальной токовой защиты (МТЗ 1 ст, МТЗ 2 ст, МТЗ 3 ст)}}
%\color{black}

\item Ступени максимальной токовой защиты (МТЗ 1 ст, МТЗ 2 ст, МТЗ 3 ст)

\begin{enumerate}[label=\arabic{section}.\arabic{subsection}.\arabic{enumi}.\arabic*, labelsep=4pt, leftmargin=0em, itemindent=65pt, parsep=0pt]

\item
Ввод каждой ступени МТЗ в работу на этапе параметрирования устройства осуществляется своим программным переключателем SGF1 <<Ввод\_функции>>. Информация о введенной ступени отображается выходным сигналом для каждой ступени (<<МТЗ 1 ст: Ввод>>, <<МТЗ 2 ст: Ввод>>, <<МТЗ 3 ст: Ввод>>).
\item
Каждая из ступеней <<МТЗ>> может быть выведена из работы оперативно путем активации сигнала оперативного вывода (<<ОВ МТЗ 1 ст>>, <<ОВ МТЗ 2 ст>>, <<ОВ МТЗ 3 ст>>), что характеризуется активным состоянием сигнала <<Оперативный вывод>> (<<МТЗ 1 ст: Оперативный вывод>>, <МТЗ 2 ст: Оперативный вывод>>, <МТЗ 3 ст: Оперативный вывод>>) на выходе функции.
\item
При активном состоянии входного сигнала <<МТЗ 1 ст на сигн>>, <<МТЗ 2 ст на сигн>>, <<МТЗ 3 ст на сигн>> действие соответствующей ступени МТЗ переводится на сигнал.
\item
Функции <<МТЗ 1 ст>>, <<МТЗ 2 ст>>, <<МТЗ 3 ст>> контролируют повышение любого фазного тока (\textsl{IA, IB, IC}) или линейного тока (\textsl{IAB, IBC, ICA}) (в зависимости от положения программного переключателя SGF1 <<Сборка\_ток\_цепей>>) выше заданной уставки, при превышении заданной уставки по току через выдержку времени таймера T1 (уставка <<Тср>>) выдается сигнал на отключение трансформатора. 
\item
Для ступени МТЗ возможны два режима работы с контролем напряжения секций НН: <<Вольтметровая блокировка>> и <<Управляющее напряжение>> (выбор режима работы осуществляется программным переключателем SGF2 <<Тип\_КПОН>>). Для выбора режима работы без контроля напряжения секций НН необходимо перевести программные переключатели SGF5 <<Реж\_КПОН1>> и SGF6 <<Реж\_КПОН2>> в состояние <<Не предусмотрено>>.
\item
В режиме <<Вольтметровая блокировка>> при превышении уставки по току для срабатывания через выдержку времени таймера T1 (уставка <<Тср>>) необходимо подтверждение от КПОН (снижение напряжения на секции НН). При неисправностях в цепях напряжения ступень МТЗ может или выводиться из работы, или действовать на отключение, в зависимости от выбранного состояния программного переключателя SGF4 <<Реж\_БНН\_КПОН>>: <<Деблокировка>> или <<Блокировка>>. 
\item
Для работы в режиме <<Управляющее напряжение>> необходимо задать две уставки по току для ступени МТЗ: <<Iср>> -- уставка по току срабатывания и <<Iср\_груб>> -- уставка по ток срабатывания грубого органа. В указанном режиме при отсутствии неисправности цепей напряжения при превышении уставки по току <<Iср>> для срабатывания необходимо подтверждение от КПОН (снижение напряжения на секции НН). При неисправностях в цепях напряжения и при превышении уставки по току грубого органа <<Iср\_груб>> через выдержку времени таймера T1 (уставка T1 <<Тср>>) выдается сигнал срабатывания. Режим работы на уставке <<Iср\_груб>> определяется состоянием программного переключателя SGF4 <<Реж\_БНН\_КПОН>>: <<Деблокировка>> или <<Блокировка>>.
\item
При необходимости возможна блокировка выбранных ступеней МТЗ при броске тока намагничивания силового трансформатора от функции  <<БНТ>>, это достигается путем перевода программного переключателя выбора режима контроля от БНТ SGF3 <<Реж\_БНТ>> в положение <<Предусмотрено>>.
\item
Для ступеней функции <<МТЗ>> предусмотрена возможность отслеживания положения СВ стороны НН. Для правильной работы в этом режиме необходимо выделить две ступени МТЗ и выставить программный переключатель SGF7 <<Контр\_СВ>> в положение <<Блокировка ступени при включенном СВ>> для ступени с ускоренной выдержкой времени и в положение <<Блокировка ступени при отключенном СВ>> для ступени с нормальной выдержкой времени. В этом случае, при изменении состояния СВ НН, автоматически вводится в работу одна из ступеней со своими уставками, учитывающими изменение режима работы сети.
\item
Алгоритм работы ступеней МТЗ выполнен в соответствии с рисунком \ref{mtzvn:img1}. В таблице \ref{mtzvn:tbl1} приведены параметры, необходимые для настройки функции <<МТЗ>>.

\small
\begin{longtable}{|>{\centering\arraybackslash}m{5.3cm}|>{\centering\arraybackslash}m{3.3cm}|>{\centering\arraybackslash}m{4.2cm}|>{\centering\arraybackslash}m{1.8cm}|>{\centering\arraybackslash}m{1cm}|}
\caption{Параметры для настройки <<МТЗ>>\hfill\vspace{-0.5\baselineskip}}\label{mtzvn:tbl1}\\ 
\hline
\rowcolor{gray!30}
Параметр (Параметр на ИЧМ) & Условное обозначение на схеме & Значение/ Диапазон & Единица измерения & Шаг \\ 
\hline
\endfirsthead
%\caption{\emph{Продолжение\hfill\vspace{-0.5\baselineskip}}} \\ % Отступ обозначения таблицы от таблицы и выравнивание надписи слева
\caption*{\hspace{3pt}\emph{Продолжение таблицы \ref{mtzvn:tbl1}\hfill\vspace{-0.5\baselineskip}}} \\ % сделано по ГОСТ 2.105 п.6.8.7
\hline
\rowcolor{gray!30}
Параметр (Параметр на ИЧМ) & Условное обозначение на схеме & Значение/ Диапазон & Единица измерения & Шаг \\ 
%\hline
\endhead
%\hline
\endfoot
%\hline
\endlastfoot
%==+t1*LLN0|T2_LVTTOC> Общие уставки
%===t1*LLN0> Общие уставки
\multicolumn{5}{|c|}{ Общие уставки } \\ \hline 
\centering Сборка токовых цепей (Сборка\_ток\_цепей) & \centering SGF1 & \centering 0 = Звезда\\1 = Треугольник & \centering -- & \centering \arraybackslash -- \\
\hline
%===t1
%==+t1*PTOC1|T2_LVTTOC> Ступень МТЗ
%===t1*PTOC1> Ступень МТЗ
\multicolumn{5}{|c|}{ Ступень МТЗ } \\ \hline 
\centering Ввод функции в работу (Ввод\_функции) & \centering SGF1 & \centering 0 = Не предусмотрено\\1 = Предусмотрено & \centering -- & \centering \arraybackslash -- \\
\hline
\centering Ток срабатывания (Iср) & \centering I$>$ & \centering 0,10 ... 30,00 & \centering о.е. & \centering \arraybackslash 0,01 \\
\hline
\centering Ток срабатывания грубого органа в режиме управляющего напряжения (Iср\_груб) & \centering Iгруб$>$ & \centering 0,10 ... 30,00 & \centering о.е. & \centering \arraybackslash 0,01 \\
\hline
\centering Тип пуска по напряжению (Тип\_КПОН) & \centering SGF2 & \centering 1 = Управляющее напряжение\\2 = Вольтметровая блокировка & \centering -- & \centering \arraybackslash -- \\
\hline
\centering Режим контроля от БНТ (Реж\_БНТ) & \centering SGF3 & \centering 0 = Не предусмотрено\\1 = Предусмотрено & \centering -- & \centering \arraybackslash -- \\
\hline
\centering Режим КПОН при неисправности ЦН (Реж\_БНН\_КПОН) & \centering SGF4 & \centering 1 = Деблокировка\\2 = Блокировка & \centering -- & \centering \arraybackslash -- \\
\hline
\centering Режим контроля от КПОН1 (Реж\_КПОН1) & \centering SGF5 & \centering 0 = Не предусмотрено\\1 = Предусмотрено & \centering -- & \centering \arraybackslash -- \\
\hline
\centering Режим контроля от КПОН2 (Реж\_КПОН2) & \centering SGF6 & \centering 0 = Не предусмотрено\\1 = Предусмотрено & \centering -- & \centering \arraybackslash -- \\
\hline
\centering Режим контроля СВ НН (Контр\_СВ) & \centering SGF7 & \centering 0 = Не предусмотрено\\1 = Блокировка ступени при включенном СВ\\2 = Блокировка ступени при отключенном СВ & \centering -- & \centering \arraybackslash -- \\
\hline
\centering Независимая выдержка времени срабатывания (Tср) & \centering T1 & \centering 0,00 ... 100,00 & \centering с & \centering \arraybackslash 0,01 \\
\hline
%===t1
\end{longtable}
\normalsize

\vspace{3mm}
\begin{figure}[!h]
\centering
\includegraphics[width=1\textwidth,height=1\textheight,keepaspectratio]{\fbpath/24. МТЗ Т2/_latex/img/mtzst.pdf}
\captionof{figure}{Функциональная схема ступеней <<МТЗ 1 ст>>, <<МТЗ 2 ст>>, <<МТЗ 3 ст>>}\label{mtzvn:img1}
\end{figure}

\end{enumerate}
%%%%%%%%%%%%%%%%%%%%%%%%%%%%%%%%%%%%%%%%%%%%%%%%%%%%%%%%% ОРГАН КПОН %%%%%%%%%%%%%%%%%%%%%%%%%%%%%%%%%%%%%%%%%%%%%%%%

%\color{uniblue}{\subsubsection{Комбинированный пусковой орган напряжения (КПОН)}}
%\color{black}

\item Комбинированный пусковой орган напряжения (КПОН)

\begin{enumerate}[label=\arabic{section}.\arabic{subsection}.\arabic{enumi}.\arabic*, labelsep=4pt, leftmargin=0em, itemindent=65pt, parsep=0pt]

\item
В составе устройства предусматриваются две функции <<КПОН1>> и <<КПОН2>> (для стороны НН1 и НН2 соответственно), ниже приводится общее описание работы функции <<КПОН>>. 
\item
Ввод функции <<КПОН>> осуществляется автоматически при активации режим пуска от внутреннего <<КПОН>> введенной в работу ступени МТЗ.
\item
Положением программного переключателя SGF1 <<Реж\_пуска>> осуществляется выбор алгоритма работы:
\begin{itemize}
\item по снижению минимального из трех линейных напряжений $U_{ab}, U_{bc}, U_{ca}$ ниже уставки $U_{\textsl{мин}}$ (состояние программного переключателя <<По Uмин>>);
\item комбинированный: по снижению минимального из трех линейных напряжений $U_{ab}, U_{bc}, U_{ca}$ ниже уставки $U_{\textsl{мин}}$ или по повышению напряжения обратной последовательности $U2$ выше уставки $U2_{\textsl{макс}}$ (состояние программного переключателя <<Комбинированный>>);
\item при приеме сигнала от внешнего КПОН <<КПОН внеш.>> (состояние программного переключателя <<Внешний>>).
\end{itemize}
\item
Функция <<КПОН>> контролирует снижение любого из линейных напряжений (\textsl{UAB, UBC, UCA}) ниже заданной уставки или, при введенном положении <<Комбинированный>> программного переключателя <<SGF1>>, повышение напряжения обратной последовательности (\textsl{U2}) выше заданной уставки или прием сигнала <<КПОН внеш.>>, в случае выполнения заданных условий выдается сигнал пуска КПОН.
\item
В качестве воздействующих величин функции <<КПОН>> использует действующие значения основной гармоники линейных напряжений и напряжение обратной последовательности. Коэффициент возврата пороговых элементов равен 1,05 для ИО по минимальному напряжению и 0,95 для ИО по максимальному напряжению.
\item
Алгоритм работы <<КПОН>> выполнен в соответствии с рисунком \ref{mtzvn:img2}. В таблице \ref{mtzvn:tbl2} приведены параметры, необходимые для настройки функции <<КПОН>>.

\vspace{3mm}
\begin{figure}[!h]
\centering
\includegraphics[width=0.9\textwidth,height=0.9\textheight,keepaspectratio]{\fbpath/24. МТЗ Т2/_latex/img/kpon.pdf}
\captionof{figure}{Функциональная схема <<КПОН>>}\label{mtzvn:img2}
\end{figure}

\small
\begin{longtable}{|>{\centering\arraybackslash}m{5.3cm}|>{\centering\arraybackslash}m{3.3cm}|>{\centering\arraybackslash}m{4.2cm}|>{\centering\arraybackslash}m{1.8cm}|>{\centering\arraybackslash}m{1cm}|}
\caption{Параметры для настройки функции <<КПОН>>\hfill\vspace{-0.5\baselineskip}}\label{mtzvn:tbl2}\\ 
\hline
\rowcolor{gray!30}
Параметр (Параметр на ИЧМ) & Условное обозначение на схеме & Значение/ Диапазон & Единица измерения & Шаг \\ 
\hline
\endfirsthead
%\caption{\emph{Продолжение\hfill\vspace{-0.5\baselineskip}}} \\ % Отступ обозначения таблицы от таблицы и выравнивание надписи слева
\caption*{\hspace{3pt}\emph{Продолжение таблицы \ref{mtzvn:tbl2}\hfill\vspace{-0.5\baselineskip}}} \\ % сделано по ГОСТ 2.105 п.6.8.7
\hline
\rowcolor{gray!30}
Параметр (Параметр на ИЧМ) & Условное обозначение на схеме & Значение/ Диапазон & Единица измерения & Шаг \\ 
%\hline
\endhead
%\hline
\endfoot
%\hline
\endlastfoot
%==+t1*PTUV1|T2_LVTTOC
%===t1*PTUV1
\centering Напряжение срабатывания (Uср) & \centering U$<$ & \centering 2,0 ... 100,0 & \centering В & \centering \arraybackslash 0,1 \\
\hline
\centering Напряжение срабатывания ОП (U2ср) & \centering U2$>$ & \centering 2,0 ... 50,0 & \centering В & \centering \arraybackslash 0,1 \\
\hline
\centering Режим пуска (Реж\_пуска) & \centering SGF1 & \centering 0 = По Uмин\\1 = Комбинированный\\2 = Внешний & \centering -- & \centering \arraybackslash -- \\
\hline
%===t1
\end{longtable}
\normalsize

\end{enumerate}
%%%%%%%%%%%%%%%%%%%%%%%%%%%%%%%%%%%%%%%%%%%%%%%%%%%%%%%%% ОРГАН БНТ %%%%%%%%%%%%%%%%%%%%%%%%%%%%%%%%%%%%%%%%%%%%%%%%

%\color{uniblue}{\subsubsection{Блокировка при броске намагничивающего тока (БНТ)}}
%\color{black}

\item Блокировка при броске намагничивающего тока (БНТ)

\begin{enumerate}[label=\arabic{section}.\arabic{subsection}.\arabic{enumi}.\arabic*, labelsep=4pt, leftmargin=0em, itemindent=65pt, parsep=0pt]

\item
Функция <<БНТ>> позволяет заблокировать защиту при броске тока намагничивания, возникающем при постановке силового трансформатора под напряжение. 
\item
Ввод функции <<БНТ>> осуществляется автоматически при активации режим контроля броска тока намагничивания введенной в работу ступени МТЗ.
\item
Условие пуска выполняется при одновременном наличии следующих условий:
\begin{itemize}
    \item отношение действующих значений фазных токов второй гармоники (\textsl{IA 2г, IB 2г, IC 2г}) к действующим значениям фазных токов основной гармоники (\textsl{IA, IB, IC}) превышает уставку <<Ih2/Ih1>> (контролируется для каждой фазы отдельно);
    \item отсутствует срабатывание порогового элемента, реагирующего на превышение максимального из действующих значений основной гармоники фазных токов (\textsl{IA, IB, IC}) над уставкой <<I\_ср>>;
    \item сформирован пуск одной из ступеней МТЗ (контроль выполняется пофазный).
\end{itemize}

\item
В функции предусмотрена возможность ввода перекрестной блокировки: при переводе программного переключателя SGF1 <<Перекрест\_блок>> в состояние <<Предусмотрено>> обнаружение броска тока намагничивания в одной из фаз блокирует срабатывание по всем фазам ступени МТЗ. Пуск функции <<БНТ>> учитывается отдельно для каждой из ступеней МТЗ при помощи соответствующего программного переключателя <<Реж\_БНТ>> в составе ступени МТЗ. При этом блокировка каждой ступени осуществляется пофазно.
\item
Алгоритм работы функции <<БНТ>> выполнен в соответствии с рисунком \ref{mtzvn:img3}. В таблице \ref{mtzvn:tbl3} приведены параметры, необходимые для настройки функции <<БНТ>>.

\vspace{3mm}
\begin{figure}[!h]
\centering
\includegraphics[width=\textwidth,height=\textheight,keepaspectratio]{\fbpath/24. МТЗ Т2/_latex/img/bnt.pdf}
\captionof{figure}{Функциональная схема БНТ}\label{mtzvn:img3}
\end{figure}

\small
\begin{longtable}{|>{\centering\arraybackslash}m{5.3cm}|>{\centering\arraybackslash}m{3.3cm}|>{\centering\arraybackslash}m{4.2cm}|>{\centering\arraybackslash}m{1.8cm}|>{\centering\arraybackslash}m{1cm}|}
\caption{Параметры для настройки функции <<БНТ>>\hfill\vspace{-0.5\baselineskip}}\label{mtzvn:tbl3}\\ 
\hline
\rowcolor{gray!30}
Параметр (Параметр на ИЧМ) & Условное обозначение на схеме & Значение/ Диапазон & Единица измерения & Шаг \\ 
\hline
\endfirsthead
%\caption{\emph{Продолжение\hfill\vspace{-0.5\baselineskip}}} \\ % Отступ обозначения таблицы от таблицы и выравнивание надписи слева
\caption*{\hspace{3pt}\emph{Продолжение таблицы \ref{mtzvn:tbl3}\hfill\vspace{-0.5\baselineskip}}} \\ % сделано по ГОСТ 2.105 п.6.8.7
\hline
\rowcolor{gray!30}
Параметр (Параметр на ИЧМ) & Условное обозначение на схеме & Значение/ Диапазон & Единица измерения & Шаг \\ 
%\hline
\endhead
%\hline
\endfoot
%\hline
\endlastfoot
%==+t1*PHAR1|T2_LVTTOC
%===t1*PHAR1
\centering Отношение тока второй гармоники к току основной гармоники (Ih2/Ih1) & \centering I2г/I1г$>$ & \centering 5 ... 100 & \centering \% & \centering \arraybackslash 1 \\
\hline
\centering Ток блокировки (Iср) & \centering I$>$ & \centering 0,1 ... 15,0 & \centering о.е. & \centering \arraybackslash 0,1 \\
\hline
\centering Перекрестная блокировка (Перекрест\_блок) & \centering SGF1 & \centering 0 = Не предусмотрено\\1 = Предусмотрено & \centering -- & \centering \arraybackslash -- \\
\hline
%===t1
\end{longtable}
\normalsize

\end{enumerate}
%%%%%%%%%%%%%%%%%%%%%%%%%%%%%%%%%%%%%%%%%%%%%%%%%%%%%%%%% ОРГАН БЛЗШ %%%%%%%%%%%%%%%%%%%%%%%%%%%%%%%%%%%%%%%%%%%%%%%%

%\color{uniblue}{\subsubsection{Блокировка логической защиты шин (БЛЗШ)}}
%\color{black}

\item Блокировка логической защиты шин (БЛЗШ)

\begin{enumerate}[label=\arabic{section}.\arabic{subsection}.\arabic{enumi}.\arabic*, labelsep=4pt, leftmargin=0em, itemindent=65pt, parsep=0pt]

\item
Функция <<БЛЗШ>> предусматривается для блокировки ЛЗШ на вводном и секционном выключателях секции, питающей силовой трансформатор. В качестве воздействующих величин используются сигналы пуска ступеней МТЗ. Состояние программного переключателя SGF1 <<БлокЛЗШ\_выбор\_ст>> определяет ступень МТЗ, от которой будет производиться блокировка ЛЗШ.
\item
Алгоритм работы <<БЛЗШ>> выполнен в соответствии с рисунком \ref{mtzvn:img4}. В таблице \ref{mtzvn:tbl4} приведены параметры, необходимые для настройки функции <<БЛЗШ>>.

\begin{figure}[!h]
\centering
\includegraphics[width=0.65\textwidth,height=0.65\textheight,keepaspectratio]{\fbpath/24. МТЗ Т2/_latex/img/blzsh.pdf}
\captionof{figure}{Функциональная схема БЛЗШ}\label{mtzvn:img4}
\vspace{3mm}
\end{figure}

\small
\begin{longtable}{|>{\centering\arraybackslash}m{5.3cm}|>{\centering\arraybackslash}m{3.3cm}|>{\centering\arraybackslash}m{4.2cm}|>{\centering\arraybackslash}m{1.8cm}|>{\centering\arraybackslash}m{1cm}|}
\caption{Параметры для настройки функции <<БЛЗШ>>\hfill\vspace{-0.5\baselineskip}}\label{mtzvn:tbl4}\\ 
\hline
\rowcolor{gray!30}
Параметр (Параметр на ИЧМ) & Условное обозначение на схеме & Значение/ Диапазон & Единица измерения & Шаг \\ 
\hline
\endfirsthead
%\caption{\emph{Продолжение\hfill\vspace{-0.5\baselineskip}}} \\ % Отступ обозначения таблицы от таблицы и выравнивание надписи слева
\caption*{\hspace{3pt}\emph{Продолжение таблицы \ref{mtzvn:tbl4}\hfill\vspace{-0.5\baselineskip}}} \\ % сделано по ГОСТ 2.105 п.6.8.7
\hline
\rowcolor{gray!30}
Параметр (Параметр на ИЧМ) & Условное обозначение на схеме & Значение/ Диапазон & Единица измерения & Шаг \\ 
%\hline
\endhead
%\hline
\endfoot
%\hline
\endlastfoot
%==+t1*RBLC1|T2_LVTTOC
%===t1*RBLC1
\centering Выбор ступени блокировки (БлокЛЗШ\_выбор\_ст) & \centering SGF1 & \centering 1 = Не предусмотрено\\2 = 1 ступень\\3 = 2 ступень\\4 = 3 ступень & \centering -- & \centering \arraybackslash -- \\
\hline
%===t1
\end{longtable}
\normalsize

\end{enumerate}

\end{enumerate}
\FloatBarrier % Все float'ы ДО этой строки будут размещены здесь.